\documentclass[11pt]{article}
\usepackage[margin=1in]{geometry}
\usepackage{amsmath, amssymb, amsthm}
\usepackage{mathtools}
\usepackage{hyperref}
\usepackage{enumitem}

\newcommand{\R}{\mathbb{R}}
\newcommand{\C}{\mathbb{C}}
\newcommand{\cG}{\mathcal{G}}
\newcommand{\cM}{\mathcal{M}}
\newcommand{\E}{\mathbb{E}}
\newcommand{\Prb}{\mathbb{P}}

\theoremstyle{plain}
\newtheorem{proposition}{Proposition}
\theoremstyle{remark}
\newtheorem{remark}{Remark}
\theoremstyle{definition}
\newtheorem{example}{Example}

\title{Notes on Model Validation in Bayesian Inverse Problems\\[6pt]
\large With Application to Disorder Inference in Majorana Nanowires}
\author{Nathan Glatt-Holtz}
\date{Last updated: \today}

\begin{document}
\maketitle
\tableofcontents

%% ===================================================================
\section{Setting and Motivation}
%% ===================================================================

\subsection{The general Bayesian inverse problem}

We consider the standard framework for Bayesian inversion. Let 
$\theta \in \Theta$ be an unknown parameter (possibly infinite-dimensional), 
$\cG: \Theta \to \R^d$ a forward map, and $\mathbf{y} \in \R^d$ observed 
data related to $\theta$ by the measurement model
\begin{equation}\label{eq:measurement}
\mathbf{y} = \cG(\theta) + \eta, \qquad \eta \sim p_\eta,
\end{equation}
where $\eta$ is observation noise with known (or partially known) distribution 
$p_\eta$. In the Gaussian case, $\eta \sim N(0, \Sigma)$ with known covariance 
$\Sigma \in \R^{d \times d}$, and the likelihood is
\begin{equation}\label{eq:likelihood}
L(\mathbf{y} \mid \theta) = (2\pi)^{-d/2}|\Sigma|^{-1/2}
\exp\Bigl(-\tfrac{1}{2}\|\Sigma^{-1/2}(\mathbf{y} - \cG(\theta))\|^2\Bigr).
\end{equation}
Given a prior distribution $\pi_0(\theta)$, the posterior is
\begin{equation}\label{eq:posterior}
\pi(\theta \mid \mathbf{y}) = \frac{L(\mathbf{y} \mid \theta)\,\pi_0(\theta)}{Z},
\qquad Z := \int_\Theta L(\mathbf{y} \mid \theta)\,\pi_0(\theta)\,d\theta,
\end{equation}
where $Z$ is the \textbf{evidence} (marginal likelihood). In practice, 
$\pi(\theta \mid \mathbf{y})$ is explored via MCMC, which produces samples 
$\theta^{(1)}, \ldots, \theta^{(K)} \sim \pi(\theta \mid \mathbf{y})$ but 
does not directly compute $Z$.

\subsection{The model criticism question}

The posterior~\eqref{eq:posterior} is the optimal answer to the question: 
``Given that the measurement model~\eqref{eq:measurement} is correct, what 
can we learn about $\theta$?'' But it does not answer the prior question: 
\emph{is the model correct?}

Model misspecification can take several forms:
\begin{enumerate}[label=(\roman*)]
\item The forward map $\cG$ does not capture the true physics (e.g., the 
  physical model is missing degrees of freedom, or uses an incorrect 
  constitutive law).
\item The noise model $p_\eta$ is wrong (e.g., the noise is non-Gaussian, 
  or the covariance $\Sigma$ is misspecified).
\item The prior $\pi_0$ is too restrictive (e.g., the true parameter lies 
  outside the support of $\pi_0$, or the smoothness assumptions are wrong).
\item The parameterization of $\theta$ is insufficient (e.g., the unknown 
  is really a function but is being approximated by too few basis coefficients).
\end{enumerate}

This document describes four complementary approaches to detecting and 
diagnosing these failures within the Bayesian framework:
posterior predictive checks (\S\ref{sec:ppc}), 
Bayesian model comparison via the evidence (\S\ref{sec:evidence}), 
hierarchical inference of hyperparameters (\S\ref{sec:hierarchical}), and 
residual analysis (\S\ref{sec:residuals}).

%% ===================================================================
\section{Posterior Predictive Checks}\label{sec:ppc}
%% ===================================================================

\subsection{The posterior predictive distribution}

Given posterior samples $\theta^{(k)} \sim \pi(\theta \mid \mathbf{y})$, the 
\textbf{posterior predictive distribution} is the distribution of hypothetical 
future data $\mathbf{y}^{\text{rep}}$ generated under the fitted model:
\begin{equation}\label{eq:ppd}
p(\mathbf{y}^{\text{rep}} \mid \mathbf{y}) 
= \int L(\mathbf{y}^{\text{rep}} \mid \theta)\,\pi(\theta \mid \mathbf{y})\,d\theta.
\end{equation}
In practice this integral is evaluated by Monte Carlo: for each posterior 
sample $\theta^{(k)}$, draw $\mathbf{y}^{(k)} \sim L(\cdot \mid \theta^{(k)})$, 
i.e.,
\begin{equation}\label{eq:ppd_sample}
\theta^{(k)} \sim \pi(\theta \mid \mathbf{y}), \qquad
\mathbf{y}^{(k)} = \cG(\theta^{(k)}) + \eta^{(k)}, \qquad
\eta^{(k)} \sim p_\eta \;\text{(fresh, independent noise)}.
\end{equation}
The collection $\{\mathbf{y}^{(1)}, \ldots, \mathbf{y}^{(K)}\}$ is a sample 
from~\eqref{eq:ppd}.

\subsection{Summary statistics and the posterior predictive p-value}

The key idea: if the model is correct, the observed data $\mathbf{y}$ should 
be a ``typical'' draw from the posterior predictive 
distribution~\eqref{eq:ppd}. To make this precise, choose a scalar 
\textbf{summary statistic} (also called a \textbf{test quantity}) 
$T: \R^d \to \R$ and compare $T(\mathbf{y})$ to the distribution of 
$T(\mathbf{y}^{(k)})$.

\begin{enumerate}[label=\textbf{Step \arabic*.}]
\item Generate $K$ replicated datasets $\mathbf{y}^{(1)}, \ldots, 
  \mathbf{y}^{(K)}$ via~\eqref{eq:ppd_sample}. (You already have 
  the posterior samples from MCMC; the only additional cost is 
  evaluating $\cG(\theta^{(k)})$, which you have already done during 
  the MCMC run, plus drawing fresh noise $\eta^{(k)}$.)
\item Compute the histogram of $\{T(\mathbf{y}^{(1)}), \ldots, 
  T(\mathbf{y}^{(K)})\}$.
\item Compute $T(\mathbf{y})$ for the actual observed data.
\item Report the \textbf{posterior predictive p-value}:
  \begin{equation}\label{eq:ppp}
  p_B := \frac{1}{K}\sum_{k=1}^K \mathbf{1}\bigl[T(\mathbf{y}^{(k)}) 
  \geq T(\mathbf{y})\bigr].
  \end{equation}
\end{enumerate}

\noindent\textbf{Interpretation.} If $p_B$ is close to 0 or 1, the observed 
data is in the tail of the predictive distribution for that statistic, which 
is evidence that the model cannot reproduce the feature captured by $T$. 
Values in $[0.1, 0.9]$ indicate no evidence of misspecification (for that 
particular $T$).

\subsection{How to choose the summary statistic $T$}

The choice of $T$ is the art of posterior predictive checking. There is no 
single ``right'' choice; different statistics probe different aspects of the 
model. The general principles are:
\begin{enumerate}[label=(\alph*)]
\item \textbf{Omnibus statistics} test overall fit. The most common is the 
  chi-squared discrepancy:
  \begin{equation}\label{eq:T_chi2}
  T_{\chi^2}(\mathbf{y}, \theta) 
  = \|\Sigma^{-1/2}(\mathbf{y} - \cG(\theta))\|^2.
  \end{equation}
  Note that this $T$ depends on both $\mathbf{y}$ and $\theta$, so the 
  p-value~\eqref{eq:ppp} becomes
  $p_B = \frac{1}{K}\sum_k \mathbf{1}[T_{\chi^2}(\mathbf{y}^{(k)}, 
  \theta^{(k)}) \geq T_{\chi^2}(\mathbf{y}, \theta^{(k)})]$.
  This catches overall lack of fit but can miss structured discrepancies.

\item \textbf{Targeted statistics} probe specific features:
  \begin{itemize}
  \item Location of a peak: $T(\mathbf{y}) = \arg\max_i y_i$.
  \item Symmetry: $T(\mathbf{y}) = \|y_{\text{left}} - y_{\text{right}}\|$ 
    for data components that should be symmetric under the model.
  \item Smoothness: $T(\mathbf{y}) = \sum_i (y_{i+1} - y_i)^2$ (total 
    variation).
  \item Tail behavior: $T(\mathbf{y}) = \max_i |y_i|$.
  \end{itemize}

\item \textbf{Use multiple statistics.} Check 3--5 different $T$'s, each 
  probing a different model assumption. If all give $p_B \in [0.1, 0.9]$, 
  there is no evidence of misspecification. If one gives an extreme $p_B$, 
  the specific $T$ that flagged it tells you \emph{what} the model is 
  getting wrong.
\end{enumerate}

\subsection{Limitations of posterior predictive checks}

\begin{enumerate}[label=(\roman*)]
\item \textbf{Double use of data.} The data $\mathbf{y}$ is used both to 
  fit the posterior and to compute $T(\mathbf{y})$. This makes the 
  p-value~\eqref{eq:ppp} conservative: it is biased toward 0.5 and can 
  fail to detect real misspecification. The severity of this bias depends 
  on the ``effective number of parameters'' in the model --- it is worst 
  when the model is highly flexible.

\item \textbf{Cross-validation variants.} To avoid the double-use problem, 
  one can hold out a subset of the data $\mathbf{y}_{\text{test}}$ and use 
  the remaining data $\mathbf{y}_{\text{train}}$ to fit the posterior. Then 
  the predictive check is performed on $\mathbf{y}_{\text{test}}$. This is 
  sharper but more expensive (multiple MCMC runs, or importance-weight 
  resampling from a single run).

\item \textbf{No formal ``rejection.''} Posterior predictive p-values are 
  not classical p-values; they do not have a uniform null distribution 
  under the model. They are best treated as continuous diagnostics, not 
  binary hypothesis tests.
\end{enumerate}


%% ===================================================================
\section{Bayesian Model Comparison via the Evidence}\label{sec:evidence}
%% ===================================================================

\subsection{The Bayes factor}

Suppose we have two candidate models $\cM_1$ and $\cM_2$, each specifying 
a forward map $\cG_i$, a noise model $p_{\eta,i}$, and a prior $\pi_{0,i}$ 
on the respective parameter spaces $\Theta_i$. The evidence for model $\cM_i$ 
is
\begin{equation}\label{eq:evidence}
Z_i := \int_{\Theta_i} L_i(\mathbf{y} \mid \theta)\,\pi_{0,i}(\theta)\,d\theta,
\end{equation}
which is the normalizing constant of the posterior~\eqref{eq:posterior} under 
model $\cM_i$. The \textbf{Bayes factor} in favor of $\cM_1$ over $\cM_2$ is
\begin{equation}\label{eq:bayes_factor}
B_{12} := \frac{Z_1}{Z_2}.
\end{equation}
If $B_{12} > 1$, the data favor $\cM_1$; if $B_{12} < 1$, they favor $\cM_2$. 
The Jeffreys scale (a rough guideline): $B_{12} > 10$ is ``strong'' evidence, 
$B_{12} > 100$ is ``decisive.''

\begin{remark}[Built-in Occam's razor]
The evidence~\eqref{eq:evidence} is \emph{not} the likelihood maximized over 
$\theta$ (that would be the maximum likelihood, which always favors more 
complex models). It is the likelihood \emph{averaged} over the prior. A model 
with more parameters spreads its prior mass over a larger volume, diluting 
the average. So the evidence automatically penalizes unnecessary complexity: 
a simpler model that fits the data adequately will have higher evidence than 
a complex model that fits slightly better but wastes prior mass on 
parameter regions that the data don't support. This is the Bayesian Occam's 
razor.
\end{remark}

\subsection{Computing the evidence: thermodynamic integration}
\label{sec:TI}

Standard MCMC produces samples from $\pi(\theta \mid \mathbf{y})$ but does 
not compute the normalizing constant $Z$. \textbf{Thermodynamic integration} 
(TI), also called \textbf{power posterior sampling}, is a method for estimating 
$\log Z$ from a sequence of MCMC runs at different ``temperatures.''

Define the \textbf{power posterior} at inverse temperature $\beta \in [0, 1]$:
\begin{equation}\label{eq:power_posterior}
\pi_\beta(\theta) := \frac{L(\mathbf{y} \mid \theta)^\beta\,\pi_0(\theta)}{Z(\beta)},
\qquad Z(\beta) := \int L(\mathbf{y} \mid \theta)^\beta\,\pi_0(\theta)\,d\theta.
\end{equation}
At $\beta = 0$, $\pi_0(\theta) = \pi_0(\theta)$ (the prior) and $Z(0) = 1$ 
(assuming the prior is normalized). At $\beta = 1$, 
$\pi_1(\theta) = \pi(\theta \mid \mathbf{y})$ (the posterior) and $Z(1) = Z$ 
(the evidence). The key identity is:

\begin{proposition}[Thermodynamic identity]
\begin{equation}\label{eq:TI_identity}
\log Z = \int_0^1 \E_{\pi_\beta}\bigl[\log L(\mathbf{y} \mid \theta)\bigr]\,d\beta.
\end{equation}
\end{proposition}

\begin{proof}
Differentiate $\log Z(\beta) = \log \int L^\beta \pi_0\,d\theta$ with respect 
to $\beta$:
\[
\frac{d}{d\beta}\log Z(\beta) 
= \frac{\int (\log L)\,L^\beta \pi_0\,d\theta}{\int L^\beta \pi_0\,d\theta}
= \E_{\pi_\beta}[\log L].
\]
Integrate both sides from $\beta = 0$ to $\beta = 1$:
\[
\log Z(1) - \log Z(0) = \int_0^1 \E_{\pi_\beta}[\log L]\,d\beta.
\]
Since $Z(0) = 1$ (normalized prior), $\log Z(0) = 0$, giving~\eqref{eq:TI_identity}.
\end{proof}

\noindent\textbf{The algorithm:}
\begin{enumerate}[label=\textbf{Step \arabic*.}]
\item Choose a grid $0 = \beta_0 < \beta_1 < \cdots < \beta_B = 1$. A 
  common choice is a geometric spacing $\beta_j = (j/B)^c$ with $c \approx 5$, 
  which clusters points near $\beta = 0$ where the integrand changes most 
  rapidly.
\item For each $\beta_j$, run MCMC targeting $\pi_{\beta_j}(\theta)$ (the 
  power posterior at temperature $\beta_j$). This is the same MCMC algorithm 
  you already use, but with the log-likelihood multiplied by $\beta_j$.
\item From the samples $\{\theta^{(k)}_j\}_{k=1}^K$ at temperature $\beta_j$, 
  estimate
  \[
  I_j := \E_{\pi_{\beta_j}}[\log L] \approx \frac{1}{K}\sum_{k=1}^K 
  \log L(\mathbf{y} \mid \theta^{(k)}_j).
  \]
\item Approximate the integral~\eqref{eq:TI_identity} by the trapezoidal rule:
  \begin{equation}\label{eq:TI_numerical}
  \log \hat{Z} = \sum_{j=0}^{B-1} \frac{\beta_{j+1} - \beta_j}{2}
  (I_j + I_{j+1}).
  \end{equation}
\end{enumerate}

\noindent\textbf{Cost:} $B$ MCMC runs (typically $B = 10$--$30$), each of 
the same length as the original posterior run. The total cost is roughly 
$B$ times the cost of a single inference.

\begin{remark}[Variance and error control]
The trapezoidal approximation~\eqref{eq:TI_numerical} has discretization 
error $O((\Delta\beta)^2)$ (from the trapezoidal rule) plus Monte Carlo 
error $O(K^{-1/2})$ at each temperature. The geometric spacing concentrates 
grid points near $\beta = 0$, where the transition from prior to posterior 
is sharpest. In practice, $B = 20$ with $K = 5000$ samples per rung is 
usually sufficient.
\end{remark}

\subsection{Other methods for computing the evidence}

\subsubsection{Bridge sampling}

Given posterior samples $\{\theta_1^{(k)}\}$ and $\{\theta_2^{(k)}\}$ from 
two models, bridge sampling estimates $Z_1/Z_2$ directly using an optimal 
``bridge function.'' The estimator (Meng \& Wong, 1996) is:
\begin{equation}\label{eq:bridge}
\frac{\hat{Z}_1}{\hat{Z}_2} 
= \frac{\sum_{k=1}^{K_2} h(\theta_2^{(k)})\,
  L_1(\mathbf{y} \mid \theta_2^{(k)})\,\pi_{0,1}(\theta_2^{(k)})}
{\sum_{k=1}^{K_1} h(\theta_1^{(k)})\,
  L_2(\mathbf{y} \mid \theta_1^{(k)})\,\pi_{0,2}(\theta_1^{(k)})},
\end{equation}
where $h$ is the optimal bridge function 
$h(\theta) \propto [K_1 \pi_1(\theta) + K_2 \pi_2(\theta)]^{-1}$. Since $h$ 
depends on $Z_1/Z_2$ (through the posteriors), the estimator is applied 
iteratively until convergence.

Bridge sampling requires \textbf{overlap} between the two posteriors: the 
samples from model~1 must have nontrivial density under model~2 and vice 
versa. This works well for nested models (where one model is a special case 
of the other) but can fail for very different models.

\subsubsection{Nested sampling}

Nested sampling (Skilling, 2006) transforms the $d$-dimensional evidence 
integral into a one-dimensional integral over the prior mass. Define the 
``prior volume'' $X(\lambda) = \Prb_{\pi_0}[L(\mathbf{y} \mid \theta) > \lambda]$. 
Then $Z = \int_0^1 L^*(X)\,dX$ where $L^*(X)$ is the inverse of $X(\lambda)$. 
The algorithm:
\begin{enumerate}[label=(\alph*)]
\item Maintain $N$ ``live points'' drawn from the prior, confined to the 
  region where $L > \lambda_{\min}$.
\item At each iteration, remove the point with the lowest likelihood, 
  replace it with a new point drawn from $\pi_0$ subject to 
  $L > \lambda_{\min}$, and increment $\lambda_{\min}$.
\item The discarded points, weighted by the prior volume they represent, 
  give an estimate of $Z$.
\end{enumerate}
Nested sampling is popular in cosmology (the \texttt{dynesty}, 
\texttt{MultiNest}, and \texttt{UltraNest} packages). It works well for 
moderate dimensions ($\lesssim 30$) but the constrained prior sampling 
becomes expensive in high dimensions.

\subsubsection{Savage--Dickey density ratio (for nested models)}
\label{sec:savage_dickey}

If $\cM_1$ is $\cM_2$ with a subset of parameters $\phi$ fixed at a 
specific value $\phi_0$ (e.g., $\cM_1$: $\phi = 0$; $\cM_2$: $\phi$ free), 
then the Bayes factor has a remarkably simple form:
\begin{equation}\label{eq:savage_dickey}
B_{12} = \frac{\pi(\phi = \phi_0 \mid \mathbf{y}, \cM_2)}{\pi_0(\phi = \phi_0 \mid \cM_2)}.
\end{equation}
That is: the ratio of the posterior density to the prior density at the 
restricted value, both evaluated under the larger model $\cM_2$.

\textbf{In words:} if the data pull the posterior away from $\phi_0$, the 
numerator is small and $B_{12} < 1$ (the data favor the larger model). If 
the posterior concentrates \emph{at} $\phi_0$, the numerator is large and 
$B_{12} > 1$ (the simpler model suffices).

\textbf{Cost:} Nearly free once you have posterior samples from $\cM_2$. 
You just need a kernel density estimate of the marginal posterior of $\phi$ 
evaluated at $\phi_0$. The practical challenge is getting a good density 
estimate, especially if $\phi$ is high-dimensional.


%% ===================================================================
\section{Hierarchical Inference of Hyperparameters}\label{sec:hierarchical}
%% ===================================================================

\subsection{The idea}

Instead of fixing model parameters (e.g., noise level $\sigma$, 
physical constants $\alpha$) at their nominal values, treat them as 
unknowns with (weakly informative) priors and infer them jointly with the 
primary unknown $\theta$. The joint posterior is
\begin{equation}\label{eq:hierarchical}
\pi(\theta, \alpha \mid \mathbf{y}) \propto 
L(\mathbf{y} \mid \theta, \alpha)\,\pi_0(\theta \mid \alpha)\,\pi_{\text{hyper}}(\alpha),
\end{equation}
where $\alpha$ collects the hyperparameters and $\pi_{\text{hyper}}(\alpha)$ 
is a hyperprior (typically vague: wide Gaussian, inverse-gamma, etc.).

\subsection{How this detects misspecification}

\textbf{The key diagnostic:} if the marginal posterior 
$\pi(\alpha \mid \mathbf{y})$ concentrates far from the nominal 
(experimentally measured or physically expected) value $\alpha^*$, the model 
is compensating for a structural deficiency by distorting $\alpha$.

More precisely: let $\alpha^*$ be the ``correct'' value of the 
hyperparameter (known from an independent measurement or physical theory). 
Run MCMC on the joint model~\eqref{eq:hierarchical} and examine the 
marginal posterior $\pi(\alpha \mid \mathbf{y})$.
\begin{itemize}
\item If $\alpha^* \in$ the 95\% credible interval of 
  $\pi(\alpha \mid \mathbf{y})$: no tension. The model is consistent with the 
  known value.
\item If $\alpha^*$ is far in the tail of $\pi(\alpha \mid \mathbf{y})$: 
  the data are forcing $\alpha$ away from its known value. The model needs 
  $\alpha \neq \alpha^*$ to fit the data, which means something else in the 
  model is wrong. This is a specific, interpretable diagnostic: the \emph{direction} 
  in which $\alpha$ is pulled tells you what physics is missing.
\end{itemize}

\subsection{Advantages over model comparison}

Hierarchical inference is a \emph{continuous} version of model comparison. 
Rather than choosing between discrete models ($\cM_1$: fix $\alpha = \alpha^*$ 
vs.\ $\cM_2$: infer $\alpha$), you let the data tell you \emph{how much} 
$\alpha$ needs to deviate from $\alpha^*$. This is often more informative 
than a binary Bayes factor, because:
\begin{itemize}
\item It gives a direction and magnitude of tension, not just a number.
\item It does not require computing the evidence $Z$ (standard MCMC suffices).
\item It naturally handles multiple hyperparameters simultaneously.
\end{itemize}

\subsection{MCMC implementation: Gibbs-within-Metropolis}\label{sec:gibbs}

The joint posterior~\eqref{eq:hierarchical} has two qualitatively different 
blocks of unknowns: the primary unknown $\theta$ (high- or 
infinite-dimensional, e.g., a function $\delta(y)$) and the hyperparameters 
$\alpha$ (low-dimensional, e.g., a handful of scalars). The natural sampling 
strategy is a \textbf{Gibbs alternation} with specialized updates for each 
block.

\subsubsection{Block 1: Update $\theta \mid \alpha, \mathbf{y}$}

Conditional on the current hyperparameter values $\alpha$, the full 
conditional for $\theta$ is
\begin{equation}\label{eq:full_cond_theta}
\pi(\theta \mid \alpha, \mathbf{y}) \propto 
L(\mathbf{y} \mid \theta, \alpha)\,\pi_0(\theta \mid \alpha).
\end{equation}
This is exactly the posterior of the original inverse problem, but with the 
likelihood and prior evaluated at the current $\alpha$. For 
infinite-dimensional $\theta$ (function-valued unknowns), this is where 
\textbf{dimension-robust MCMC algorithms} are essential:
\begin{itemize}
\item \textbf{pCN} (preconditioned Crank--Nicolson): proposals of the form 
  $\theta' = \sqrt{1 - \beta^2}\,\theta + \beta\,\xi$, $\xi \sim \pi_0$, 
  which are reversible with respect to the prior and whose acceptance rate 
  is independent of the discretization dimension.
\item \textbf{$\infty$-HMC} (infinite-dimensional Hamiltonian Monte Carlo): 
  Hamiltonian dynamics on the Cameron--Martin space of the prior, giving 
  gradient-informed proposals that are also dimension-robust.
\item \textbf{mpCN and MESS} (multiproposal variants): generate multiple 
  proposals per step for improved parallel efficiency and mixing in 
  multimodal posteriors.
\end{itemize}
These serve as the \textbf{Metropolis-within-Gibbs} step for the $\theta$ 
block. The key observation is that conditioning on $\alpha$ does not change 
the structure of this update at all --- it only changes the numerical 
values plugged into the forward map $\cG(\theta; \alpha)$ and possibly the 
prior covariance operator $\mathcal{C}(\alpha)$. Existing MCMC code for the 
primary inference needs only to accept $\alpha$ as a parameter input rather 
than a hardcoded constant.

\subsubsection{Block 2: Update $\alpha \mid \theta, \mathbf{y}$}

Conditional on the current $\theta$, the full conditional for $\alpha$ is 
low-dimensional:
\begin{equation}\label{eq:full_cond_alpha}
\pi(\alpha \mid \theta, \mathbf{y}) \propto 
L(\mathbf{y} \mid \theta, \alpha)\,\pi_0(\theta \mid \alpha)\,
\pi_{\text{hyper}}(\alpha).
\end{equation}
Since $\alpha$ is typically 3--5 scalars, a simple Metropolis step (random 
walk, or adaptive Metropolis, or MALA) suffices. In some cases, conjugacy 
gives a direct Gibbs draw:

\begin{example}[Noise variance $\sigma^2$]
If $\Sigma = \sigma^2 I_d$ and $\sigma^2$ has an inverse-gamma prior 
$\sigma^2 \sim \text{IG}(a_0, b_0)$, then the full 
conditional~\eqref{eq:full_cond_alpha} is also inverse-gamma:
\[
\sigma^2 \mid \theta, \mathbf{y} \;\sim\; 
\text{IG}\Bigl(a_0 + \tfrac{d}{2},\; 
b_0 + \tfrac{1}{2}\|\mathbf{y} - \cG(\theta)\|^2\Bigr).
\]
This is a direct draw --- no Metropolis step needed, no forward model 
evaluation required (the residual $\mathbf{y} - \cG(\theta)$ is already 
available from Block~1).
\end{example}

\begin{example}[Physical parameters entering $\cG$ nonlinearly]
For hyperparameters like $\Delta$, $E_z$, or $\alpha_R$ that enter the 
forward map $\cG$ through the BdG Hamiltonian, there is no conjugacy. 
The full conditional~\eqref{eq:full_cond_alpha} requires evaluating 
$\cG(\theta; \alpha')$ at the proposed value $\alpha'$, which is a full 
forward model solve. However, since this is a 1D (or at most 3--4D) 
Metropolis update, even a simple random walk mixes quickly.
\end{example}

\subsubsection{The full Gibbs iteration}

One iteration of the hierarchical sampler is:
\begin{equation}\label{eq:gibbs_iteration}
\boxed{
\begin{aligned}
&\text{1.}\quad \theta^{(k+1)} \sim \pi(\theta \mid \alpha^{(k)}, \mathbf{y})
  &&\text{(pCN/mpCN/MESS --- one or several steps)}, \\
&\text{2.}\quad \alpha^{(k+1)} \sim \pi(\alpha \mid \theta^{(k+1)}, \mathbf{y})
  &&\text{(Metropolis or Gibbs for each component)}.
\end{aligned}
}
\end{equation}

\subsubsection{Cost analysis}

\textbf{Block~1} requires at least one evaluation of 
$\cG(\theta; \alpha^{(k)})$ per step (this is the dominant cost and is 
the same as in the non-hierarchical sampler).

\textbf{Block~2} cost depends on where $\alpha$ enters:
\begin{itemize}
\item If $\alpha$ enters \emph{only through the likelihood} and not 
  through $\cG$ (e.g., the noise variance $\sigma^2$ scales the 
  residual), then Block~2 is essentially \textbf{free}: no forward model 
  evaluation is needed, just a re-evaluation of the likelihood at the 
  current residual with the new $\alpha$.
\item If $\alpha$ enters \emph{through the Hamiltonian} (e.g., $\Delta$, 
  $E_z$, $\alpha_R$), then accepting a proposed $\alpha'$ in Block~2 
  requires a full forward model evaluation $\cG(\theta^{(k+1)}; \alpha')$. 
  This is one additional solve per accepted $\alpha$-proposal.
\end{itemize}

\textbf{Practical trick:} perform several $\theta$-updates (Block~1 
steps) per $\alpha$-update (Block~2 step). The $\theta$-updates are 
already paying for forward model evaluations, and the hyperparameters 
$\alpha$ typically mix faster (being low-dimensional) so they do not 
need to be updated at every iteration. A ratio of 5--10 $\theta$-steps 
per $\alpha$-step is a reasonable starting point.

\textbf{Total overhead} relative to the non-hierarchical sampler: if 
$\alpha$ only includes the noise variance (conjugate update), the 
overhead is negligible. If $\alpha$ includes physical parameters entering 
$\cG$, the overhead is roughly the fraction of Block~2 proposals that are 
accepted times one additional forward solve per acceptance --- typically 
a 10--30\% increase in total cost.


%% ===================================================================
\section{Residual Analysis}\label{sec:residuals}
%% ===================================================================

\subsection{The Bayesian residual}

Given a point estimate $\hat\theta$ (e.g., the posterior mean 
$\hat\theta = \E_\pi[\theta]$ or the MAP estimate 
$\hat\theta = \arg\max_\theta \pi(\theta \mid \mathbf{y})$), define the 
\textbf{residual vector}
\begin{equation}\label{eq:residual}
\mathbf{r}(\hat\theta) := \mathbf{y} - \cG(\hat\theta) \in \R^d.
\end{equation}
Under the model~\eqref{eq:measurement}, if $\hat\theta$ were exactly equal 
to the true parameter $\theta^*$, then 
$\mathbf{r}(\theta^*) = \eta \sim N(0, \Sigma)$.

\subsection{What to check}

Standard regression diagnostics apply:
\begin{enumerate}[label=(\alph*)]
\item \textbf{Normality.} Are the standardized residuals 
  $\tilde{r}_i = [\Sigma^{-1/2}\mathbf{r}]_i$ approximately $N(0, 1)$? 
  Check with a QQ-plot or a Shapiro--Wilk test. Departures suggest 
  non-Gaussian noise or model misspecification.
\item \textbf{Homoscedasticity.} Plot $|\tilde{r}_i|$ against the 
  predicted values $[\cG(\hat\theta)]_i$ or against an experimental 
  variable (e.g., bias voltage $V$). Patterns (e.g., larger residuals 
  at certain voltages) indicate that the noise model $\Sigma$ is wrong 
  or that the forward map $\cG$ systematically misses features at those 
  parameter values.
\item \textbf{Independence.} If the data have a natural ordering (e.g., 
  indexed by voltage), check for autocorrelation in the residuals. 
  Positive autocorrelation suggests the model is missing smooth structure 
  in the data.
\item \textbf{Omnibus chi-squared.} Under the model, 
  $\chi^2 := \|\Sigma^{-1/2}\mathbf{r}\|^2 \sim \chi^2_d$ 
  (approximately, for $\hat\theta$ close to $\theta^*$). If 
  $\chi^2 \gg d$, the model does not fit the data.
\end{enumerate}

\subsection{When residual analysis is and is not rigorous}

\textbf{The honest statement:} the distributional claim 
$\mathbf{r}(\hat\theta) \sim N(0, \Sigma)$ is only valid if:
\begin{enumerate}[label=(\roman*)]
\item The model is correctly specified (the true DGP is~\eqref{eq:measurement} 
  with $\eta \sim N(0, \Sigma)$).
\item The estimator $\hat\theta$ is close to the true $\theta^*$ (i.e., the 
  posterior has concentrated near the truth --- the Bayesian consistency 
  condition).
\end{enumerate}
If the posterior has not concentrated --- for example because the data are 
insufficient to identify $\theta^*$, or because the parameter space is 
too high-dimensional --- then $\hat\theta$ may be far from $\theta^*$ and 
the residuals $\mathbf{r}(\hat\theta)$ will reflect both the noise $\eta$ 
and the estimation error $\cG(\theta^*) - \cG(\hat\theta)$. In this case 
the residuals are not guaranteed to be Gaussian even if the model is correct.

\textbf{The practical value:} despite the lack of formal guarantees, residual 
analysis remains the fastest and most intuitive diagnostic. The logic is 
asymmetric:
\begin{itemize}
\item If the residuals look terrible (structured patterns, heavy tails, 
  heteroscedasticity), the model is almost certainly wrong.
\item If the residuals look clean, the model \emph{might} be right --- but 
  a wrong model can produce innocent-looking residuals if the misspecification 
  lies in a direction the data cannot probe.
\end{itemize}
This is the same situation as in classical regression: residual analysis 
catches many forms of misspecification but is not omniscient.

\begin{remark}[Connection to posterior predictive checks]
The chi-squared residual $\chi^2 = \|\Sigma^{-1/2}\mathbf{r}(\hat\theta)\|^2$ 
at the point estimate is closely related to the omnibus posterior predictive 
check with $T = T_{\chi^2}$ from~\eqref{eq:T_chi2}. The difference is that 
the posterior predictive check averages over posterior uncertainty in $\theta$, 
while the residual check conditions on a single point estimate. The posterior 
predictive version is more Bayesian but more expensive; the residual version 
is quick and dirty but can be sharpened by examining the residuals at 
\emph{multiple} posterior samples rather than just the point estimate.
\end{remark}


%% ===================================================================
\section{Recommendations for the Majorana Nanowire Problem}\label{sec:recs}
%% ===================================================================

\subsection{Recall of the problem setup}

We observe a conductance map 
$\mathbf{G} = \{G_{LL}^i, G_{LR}^i, G_{RL}^i, G_{RR}^i\}_{i=1}^n$ from 
a proximitized semiconductor nanowire, where the data are modeled as
\begin{equation}\label{eq:our_model}
\mathbf{G} = \cG(\delta^*) + \eta, \qquad \eta \sim N(0, \Sigma),
\end{equation}
with $\delta^*: [0, L] \to \R$ the true disorder profile and $\cG$ the 
forward map:
\[
\cG: \delta(\cdot) 
\;\xrightarrow{H_{\mathrm{BdG}}(\delta)}\; 
\;\xrightarrow{S(E;\delta)}\;
\;\xrightarrow{\text{Landauer--B\"uttiker}}\;
\mathbf{G}(\delta).
\]
The BdG Hamiltonian is 
$H_{\mathrm{BdG}} = \tau_z[-\frac{\hbar^2}{2m^*}\partial_y^2\sigma_0 
- \mu(y)\sigma_0 - i\hbar\alpha_R\sigma_x\partial_y + V_g(y)\sigma_0] 
- E_z\tau_0\sigma_z + \Delta\tau_x\sigma_0$,
where $\mu(y) = \mu_0 + \delta(y)$ and the known parameters are 
$m^*, \alpha_R, \Delta, E_z, V_g, \mu_0$. The posterior is
\[
\pi(\delta \mid \mathbf{G}) \propto 
\exp\bigl(-\tfrac{1}{2}\|\Sigma^{-1/2}(\mathbf{G} - \cG(\delta))\|^2\bigr)
\,\pi_0(\delta).
\]

\subsection{Specific model validation recommendations}

\subsubsection{Posterior predictive checks}

Generate replicated conductance maps 
$\mathbf{G}^{(k)} = \cG(\delta^{(k)}) + \eta^{(k)}$ from posterior 
samples. Recommended summary statistics:
\begin{enumerate}[label=(\roman*)]
\item \textbf{Omnibus fit:} 
  $T_1 = \|\Sigma^{-1/2}(\mathbf{G} - \cG(\delta))\|^2$ 
  (chi-squared at each posterior sample).
\item \textbf{Left--right asymmetry:} 
  $T_2 = \|G_{LL}(V=0) - G_{RR}(V=0)\|$. A true MZM pair produces 
  quantized ZBCPs at both ends; if the model is correct, the posterior 
  predictive distribution of $T_2$ should be consistent with the observed 
  asymmetry.
\item \textbf{Nonlocal conductance magnitude:} 
  $T_3 = \max_V |G_{RL}(V)|$. This probes whether the model captures 
  bulk transport correctly.
\item \textbf{ZBCP width:} 
  $T_4 = $ full width at half maximum of $G_{LL}(V)$ near $V = 0$. 
  Sensitive to temperature broadening and tunnel barrier modeling.
\end{enumerate}
\textbf{Cost:} essentially free once the MCMC run is complete, since 
$\cG(\delta^{(k)})$ has already been evaluated during sampling.

\subsubsection{Bayesian model comparison}

The most natural model comparisons for this problem:

\medskip
\noindent\textbf{(A) Disorder only vs.\ disorder $+$ unknown gate potential.}
\begin{itemize}
\item $\cM_1$: Infer $\delta(y)$ with $V_g(y) = V_g^{\text{nominal}}(y)$ 
  fixed at the experimentally specified gate profile.
\item $\cM_2$: Infer both $\delta(y)$ and $V_g(y)$ jointly.
\end{itemize}
If $\cM_2$ has decisively higher evidence, the nominal gate profile is 
inadequate and the conductance data contain information about gate 
potential variations that the disorder alone cannot explain.

\medskip
\noindent\textbf{(B) Single-band vs.\ multi-band BdG.}
\begin{itemize}
\item $\cM_1$: The current single-band model.
\item $\cM_2$: A multi-band extension with additional sub-bands in the 
  nanowire.
\end{itemize}
More expensive to implement but directly tests the most significant 
physical simplification in the model.

\medskip
\noindent\textbf{Recommended method:} Thermodynamic integration 
(\S\ref{sec:TI}), since MCMC infrastructure is already in place. For the 
parametric ($\sim 3$ parameter) version of the problem, nested sampling 
(via \texttt{dynesty}) is also practical and gives the evidence as a 
direct output.

\subsubsection{Hierarchical inference}

Promote the following nominal parameters to inferred quantities with 
weakly informative priors centered at their experimental values:
\begin{enumerate}[label=(\roman*)]
\item \textbf{Superconducting gap $\Delta$:} if the posterior on $\Delta$ 
  is pulled away from the experimentally measured value, this indicates 
  that the proximity effect is not spatially uniform or that the 
  effective gap differs from the bulk value. One could then extend the 
  model to infer a spatially varying $\Delta(y)$.
\item \textbf{Zeeman energy $E_z$:} should be well-determined by the 
  applied magnetic field, but if the effective $g$-factor varies along 
  the wire, the posterior on $E_z$ will be pulled.
\item \textbf{Noise variance $\sigma^2$} (or the full $\Sigma$): if the 
  posterior on $\sigma^2$ is much larger than the experimentally estimated 
  noise, the model is under-fitting the data (the extra variance 
  compensates for systematic model error).
\item \textbf{Spin-orbit coupling $\alpha_R$:} should be determined by 
  the material, but effective values can differ from bulk measurements 
  in confined geometries.
\end{enumerate}

\textbf{Cost:} Modest --- enlarges the parameter space by a few dimensions 
(one per hyperparameter). The MCMC mixing may slow down slightly but the 
per-sample cost of $\cG$ is unchanged.

\subsubsection{Residual analysis}

At the posterior mean $\hat\delta = \E_\pi[\delta]$, compute the 
standardized residuals 
$\tilde{r}_i = [\Sigma^{-1/2}(\mathbf{G} - \cG(\hat\delta))]_i$ and 
check:
\begin{enumerate}[label=(\roman*)]
\item \textbf{QQ-plot} of $\{\tilde{r}_i\}$ against $N(0,1)$ quantiles.
\item \textbf{Residuals vs.\ bias voltage $V$:} plot $\tilde{r}_i$ 
  against $V_i$. Systematic structure (e.g., residuals consistently 
  positive near $V = 0$ or at the gap edge $eV \approx \Delta$) 
  indicates the model misses features at specific energies.
\item \textbf{Residuals vs.\ Zeeman field $V_Z$:} plot $\tilde{r}_i$ 
  against $V_{Z,j}$. Structure near the topological phase transition 
  ($V_Z \approx \sqrt{\mu_0^2 + \Delta^2}$) would indicate the model 
  struggles precisely where the physics is most interesting.
\item \textbf{Residuals by component:} compare the residual distributions 
  for $G_{LL}$, $G_{LR}$, $G_{RL}$, $G_{RR}$ separately. If one 
  component has systematically larger residuals, the model may be 
  misspecifying the physics at one end of the wire or in the bulk.
\item \textbf{Chi-squared test:} 
  $\chi^2 = \sum_i \tilde{r}_i^2 \approx \chi^2_d$. Report $\chi^2/d$ 
  (should be $\approx 1$).
\end{enumerate}

\textbf{Cost:} Negligible (one forward model evaluation at the posterior mean).

\subsection{Recommended workflow}

In order of increasing effort:
\begin{enumerate}[label=\textbf{\arabic*.}]
\item \textbf{Residual analysis} (free, do always): inspect the 
  residuals at the posterior mean for structure. Takes minutes.
\item \textbf{Posterior predictive checks} (cheap, do always): 
  generate $K$ replicated datasets from the posterior, check 3--5 
  summary statistics. Requires no additional forward model evaluations 
  beyond what MCMC already computed.
\item \textbf{Hierarchical inference} (moderate cost, do for key 
  parameters): promote $\Delta$, $\sigma^2$, and possibly $E_z$ and 
  $\alpha_R$ to inferred quantities. Run one additional MCMC chain with 
  the expanded parameter space. Examine the marginal posteriors for 
  tension with known values.
\item \textbf{Bayesian model comparison} (expensive, do for key 
  model choices): compute the evidence via thermodynamic integration 
  for the disorder-only model vs.\ the disorder-plus-gate-potential 
  model. Requires $\sim 20$ additional MCMC runs.
\end{enumerate}


\end{document}
